\documentclass[../../main.tex]{subfiles}
\begin{document}

{\bf [解答]}

まずは一般的な状況を考える．
Aがグー，チョキ，パーを出す確率を $p, q, r$，
Bがグー，チョキ，パーを出す確率を $a, b, c$，
とする．これらは排反だから
\begin{align}
  \begin{cases}
    p + q + r = 1 \\
    a + b + c = 1
  \end{cases} \label{eq:1}
\end{align}
である．

題意で与えられたじゃんけんの手の組み合わせと歩数を表にまとめると下のようになる．

\begin{table}[htb]
  \centering
  \caption{じゃんけんの手の組み合わせと歩数}
  \begin{tabular}{|c|c|c|c|}
    \hline
    A \textbackslash{} B & グー & チョキ & パー \\ \hline
    グー & $\times$ & Bが3歩 & Aが6歩 \\ \hline
    チョキ & Aが3歩 & $\times$ & Bが5歩 \\ \hline
    パー & Bが6歩 & Aが5歩 & $\times$ \\ \hline
  \end{tabular}
\end{table}

A，Bの一回のじゃんけんでの歩数を$X_A,X_B$，その期待値を $E_A, E_B$ とすると，上表から
\begin{align}
  \begin{cases}
    E_A = E(X_A) = r \cdot a \cdot 6 + p \cdot b \cdot 3 + q \cdot c \cdot 5 \\
    E_B = E(X_B) = q \cdot a \cdot 3 + r \cdot b \cdot 5 + p \cdot c \cdot 6
  \end{cases} \label{eq:2}
\end{align}
である．Aの進む歩数からBの進む歩数を引いた値の期待値$E$は，期待値の性質から
\begin{align}
  E = E(X_A-X_B) = E(X_A) -E(X_B) = E_A-E_B
\end{align}
だから，\cref{eq:2}を代入して
\begin{align}
  E
  &= (6r - 3q) a + (3p - 5r) b + (5q - 6p) c \label{eq:3} \\
  &= (3b - 6c) p + (5c - 3a) q + (6a - 5b) r \label{eq:4}
\end{align}
となる．

\paragraph{(1)}

Bがグーチョキパーを出す確率がすべて等しい，つまり$a=b=c=\dfrac{1}{3}$の時，
\cref{eq:4}に代入して
\begin{align}
  E = -p + \dfrac{2}{3} q + \dfrac{1}{3} r \le \dfrac{2}{3}
\end{align}
である．等号成立は$p=r=0, q=1$の時である．
よって$E$が最大になるのは
\begin{align}
  (p, q, r) = (0, 1, 0)
\end{align}
の時である．

\paragraph{(2)}

\cref{eq:3}が任意の$(a,b,c)$で成立することが必要である．
\cref{eq:1}にも注意して，$(a,b,c) = (1,0,0), (0,1,0), (0,0,1)$ での成立が必要だから
\begin{align}
  &
  \begin{cases}
    6r - 3q \ge 0 \\
    3p - 5r \ge 0 \\
    5q - 6p \ge 0
  \end{cases}
  \quad \therefore
  & r \ge \dfrac{1}{2} q \ge \dfrac{3}{5} p \ge r \label{eq:5}
\end{align}
である．\cref{eq:5}が成り立つためには全ての等号が成立することが必要で
\begin{align}
  r = \dfrac{1}{2} q = \dfrac{3}{5} p
\end{align}
が必要である．
\cref{eq:1}に代入して成立すると
\begin{align}
  (p, q, r) = \Bigl( \dfrac{5}{14}, \dfrac{3}{7}, \dfrac{3}{14} \Bigr) \label{eq:6}
\end{align}
である．

逆にこの時，\cref{eq:3}は
\begin{align}
  E = 0
\end{align}
となり条件を満たす．
よって求める$(p,q,r)$の条件は
\begin{align}
  (p, q, r) = \Bigl( \dfrac{5}{14}, \dfrac{3}{7}, \dfrac{3}{14} \Bigr)
\end{align}
である．

\end{document}
